Microgrids (MG) are small-scale power subsystems that incorporate distributed generation, energy storage, local loads, system control, and coordination. Energy Management Systems (EMSs) play a crucial role in fortifying the resilience of microgrids, ensuring their stability and continued operation even when energy supply is constrained. An EMS enables centralized monitoring and control of energy consumption across various building systems. It is designed to collect, store, and analyze power consumption data from residential and industrial appliances.
\vspace{0.5cm}
To achieve a fully automated and stable grid with plug-and-play capabilities, an agent-based EMS is proposed for use in a DC microgrid. This dissertation introduces an energy management system utilizing a system of systems (SoSs) approach.  A model and simulation of an energy management system for a PV/Composite Storage DC microgrid are developed in MATLAB. The primary aim of this dissertation is to enhance the performance of the DC Microgrid by integrating an existing grid with an energy management system at the residential level. The control system's objective is to ensure prolonged operation of loads by maintaining the local battery's charge until the PV modules cease power generation. Four scenarios of varying load demand were tested. The results obtained from both scenarios demonstrate the effectiveness of this control system design approach.
\vspace{0.5cm}
In addition to digital simulations, a physical representation of the DC Microgrid (DC MG) has been constructed to validate the simulations in real-life scenarios. This physical system allows for hands-on analysis and experimentation, facilitating a deeper understanding of the data generated during the simulation phase. The combination of digital simulations and physical experimentation ensures a comprehensive and reliable assessment of the DC microgrid's behaviour and performance under various conditions.
\vspace{0.5cm}
Results indicate that non-EMS integrated DC microgrids can fulfill energy demands more directly, while EMS-integrated systems prioritise battery charging, resulting in a 10\% improvement in the state of charge of local storage. However, this feature may prevent the system from fully meeting immediate load demand. During periods of peak load, both systems experience similar SoC declines, but the EMS's energy-saving features extend battery life by reducing supply capacity when the SoC falls below a critical threshold. Despite some limitations in meeting high demand due to restricted battery capacity, the EMS enhances battery lifespan by preventing excessive discharge and improves overall grid stability. These findings suggest that a strategic balance between local and farm storage usage and grid reliance can be achieved without a central control element by using a system-of-systems approach.